\chapter{Courbes}\label{ch:b2:curves}

La géométrie devient maintenant \omterm{def:b2:diffcalc:differential}{différentielle}. Une courbe est un
point qui se déplace dans l'espace~; le calcul différentiel nous
donne sa vitesse et son accélération, et la géométrie demande ce
qui est intrinsèque --- indépendant de la vitesse à laquelle on
parcourt la trajectoire. Les réponses sont la \emph{\omterm{def:b2:curves:length}{longueur
d'arc}}, qui mesure la trajectoire elle-même, et la
\emph{\omterm{thm:b2:curves:frenet2d}{courbure}}, qui mesure comment elle se courbe. En dimension
$3$, un second invariant, la \emph{\omterm{thm:b2:curves:frenet3d}{torsion}}, mesure comment la
courbe se tord hors de son plan. L'outil qui comptabilise tout
cela est le \emph{repère de Frenet} mobile.

\section{Arcs paramétrés}

\begin{definition}[Arc paramétré]\label{def:b2:curves:arc}
Un \emph{arc paramétré}\index{arc paramétré} de classe
$\mathcal{C}^k$ ($k \geq 1$) est
une application $\gamma \colon I \to \R^n$ de classe
$\mathcal{C}^k$ sur un
intervalle $I$. Un point $\gamma(t)$ est \emph{régulier} si
$\gamma'(t) \neq 0$, et l'arc est \emph{régulier} si tous ses
points le sont. La droite passant par $\gamma(t)$ et dirigée par
$\gamma'(t)$ est la \emph{tangente}\index{tangente} en un point
régulier.
\end{definition}

\begin{definition}[Changement de paramètre]\label{def:b2:curves:reparam}
Un \emph{changement de paramètre}\index{changement de paramètre}
de classe $\mathcal{C}^k$ est un
$\mathcal{C}^k$-difféomorphisme $\theta \colon J \to I$ entre
intervalles ($\theta' \neq 0$ partout). Les arcs $\gamma$ et
$\gamma \circ \theta$ sont dits \emph{équivalents}~; un
\emph{arc géométrique}\index{arc géométrique} (ou courbe) est une
classe d'équivalence. Les notions
invariantes par changement de paramètre --- la trajectoire, la
\omterm{def:b2:curves:arc}{tangente}, la \omterm{def:b2:curves:length}{longueur d'arc}, la \omterm{thm:b2:curves:frenet2d}{courbure} --- sont dites
\emph{géométriques}.
\end{definition}

\begin{remark}
La trajectoire seule ne détermine pas l'\omterm{def:b2:curves:reparam}{arc géométrique}~: les
paramétrages $t \mapsto (\cos t, \sin t)$ sur $[0, 2\pi]$ et sur
$[0, 4\pi]$ ont la même image mais parcourent le cercle une fois
et deux fois. Un \omterm{def:b2:curves:reparam}{arc géométrique} retient la multiplicité et
l'orientation du parcours, non sa vitesse.
\end{remark}

\begin{example}[La vitesse ne change rien de
géométrique]\label{ex:b2:curves:speedinvariance}
Paramétrons le cercle unité par
\[
\gamma(t) = (\cos t^2,\ \sin t^2),
\qquad t \in \intcc0{\sqrt{2\pi}} .
\]
La vitesse $\norm{\gamma'(t)} = 2t$ croît linéairement, et
pourtant
\[
L = \int_0^{\sqrt{2\pi}}2t\,\dd t = 2\pi ,
\]
la même \omterm{def:b2:curves:length}{longueur} qu'à vitesse constante --- comme le
\cref{thm:b2:curves:lengthinv} le promet, via le \omterm{def:b2:curves:reparam}{changement de
paramètre} $t \mapsto t^2$. La \omterm{def:b2:curves:arc}{tangente}, la \omterm{thm:b2:curves:frenet2d}{courbure}
calculée à partir de la \cref{prop:b2:curves:kappaformula}, et
toute autre grandeur géométrique coïncident également~; seul $t = 0$
mérite un regard, où $\gamma'(0) = 0$ rend \emph{ce
paramétrage} irrégulier bien que la trajectoire soit un
cercle parfait. Les énoncés géométriques tolèrent mal les
mauvais paramétrages~: reparamétrer d'abord, conclure
ensuite.
\end{example}

\begin{example}\label{ex:b2:curves:cusp}
L'arc $\gamma(t) = (t^2, t^3)$ est $\mathcal{C}^\infty$ mais non
régulier~: $\gamma'(0) = (0, 0)$. Sa trajectoire, la parabole
semi-cubique $y^2 = x^3$, présente un \emph{point de
rebroussement} à l'origine~: la régularité du paramétrage
n'empêche pas une singularité géométrique là où la vitesse
s'annule. C'est pourquoi l'hypothèse de régularité
$\gamma' \neq 0$ n'est pas cosmétique.
\end{example}

\section{Longueur d'arc}

\begin{definition}[Longueur d'arc]\label{def:b2:curves:length}
Soit $\gamma \colon [a, b] \to \R^n$ un arc $\mathcal{C}^1$. Sa
\emph{longueur}\index{longueur d'arc} est
\[
L(\gamma) = \int_a^b \norm{\gamma'(t)}\, \dd t ,
\]
où $\norm{\cdot}$ est la \omterm{def:b2:nvs:norm}{norme} euclidienne. L'\emph{abscisse
curviligne}\index{abscisse curviligne} d'origine $t_0$ est $s(t) = \int_{t_0}^t
\norm{\gamma'(u)}\,\dd u$.
\end{definition}

\begin{theorem}[La longueur est géométrique~; caractérisation
polygonale]\label{thm:b2:curves:lengthinv}
\leavevmode
\begin{enumerate}
\item Si $\theta \colon [c, d] \to [a, b]$ est un \omterm{def:b2:curves:reparam}{changement de
paramètre} $\mathcal{C}^1$, alors $L(\gamma \circ \theta) = L(\gamma)$.
\item $L(\gamma)$ est la borne supérieure des \omterm{def:b2:curves:length}{longueurs} des
polygones inscrits~:
\[
L(\gamma)
= \sup\Bigl\{\, \sum_{i=1}^{m} \norm{\gamma(t_i) -
\gamma(t_{i-1})} \;:\; a = t_0 < t_1 < \dots < t_m = b \,\Bigr\}.
\]
\end{enumerate}
\end{theorem}

\begin{proof}
\emph{1.} Par le changement de variable $t = \theta(u)$ (volume
de première année, valide puisque $\theta$ est $\mathcal{C}^1$
monotone),
\[
\int_c^d \norm{(\gamma\circ\theta)'(u)}\,\dd u
= \int_c^d \norm{\gamma'(\theta(u))}\,\abs{\theta'(u)}\,\dd u
= \int_a^b \norm{\gamma'(t)}\,\dd t ,
\]
où l'on a utilisé $(\gamma\circ\theta)' = \theta'\cdot
(\gamma'\circ\theta)$ et, si $\theta$ est décroissante, le signe de
$\theta'$ est absorbé par l'inversion des bornes.

\emph{2.} Pour toute subdivision, $\gamma(t_i) - \gamma(t_{i-1}) =
\int_{t_{i-1}}^{t_i}\gamma'(t)\,\dd t$, donc par l'inégalité
triangulaire pour les intégrales $\norm{\gamma(t_i) - \gamma(t_{i-1})}
\leq \int_{t_{i-1}}^{t_i}\norm{\gamma'}$~: tout polygone est plus
court que $L(\gamma)$, d'où $\sup \leq L(\gamma)$.

Pour l'inégalité inverse, soit $\varepsilon > 0$. Comme $\gamma'$
est \omterm{def:b2:metric:continuity}{continue} sur le \omterm{def:b2:metric:compact}{compact} $[a,b]$, elle est \omterm{def:b2:funcseq:def}{uniformément}
\omterm{def:b2:metric:continuity}{continue}~: il existe $\delta > 0$ tel que $\norm{\gamma'(t) -
\gamma'(u)} \leq \varepsilon$ dès que $\abs{t - u} \leq \delta$.
Prenons une subdivision de pas $\leq \delta$. Sur chaque morceau,
pour $t \in [t_{i-1}, t_i]$,
\[
\gamma(t_i) - \gamma(t_{i-1})
= \int_{t_{i-1}}^{t_i}\gamma'(t)\,\dd t
= (t_i - t_{i-1})\,\gamma'(t_{i-1}) + R_i,
\qquad
\norm{R_i} \leq \varepsilon\,(t_i - t_{i-1}),
\]
puisque $R_i = \int_{t_{i-1}}^{t_i}(\gamma'(t) -
\gamma'(t_{i-1}))\,\dd t$. Donc
\[
\norm{\gamma(t_i) - \gamma(t_{i-1})}
\geq (t_i - t_{i-1})\norm{\gamma'(t_{i-1})}
- \varepsilon (t_i - t_{i-1}) .
\]
En sommant, et en comparant $\sum (t_i -
t_{i-1})\norm{\gamma'(t_{i-1})}$ avec $\int_a^b\norm{\gamma'}$ (une
somme de Riemann de la fonction \omterm{def:b2:metric:continuity}{continue} $\norm{\gamma'}$, à
$\varepsilon(b - a)$ près de l'intégrale pour $\delta$ assez petit,
de nouveau par continuité uniforme), on obtient un polygone de
\omterm{def:b2:curves:length}{longueur} $\geq L(\gamma) - 2\varepsilon(b - a)$. En faisant tendre
$\varepsilon \to 0$, on prouve l'assertion.
\end{proof}

\begin{example}[Archimède et les polygones
inscrits]\label{ex:b2:curves:archimedes}
Pour le cercle unité, le polygone régulier à $n$ côtés inscrit a
pour \omterm{def:b2:curves:length}{longueur} $L_n = 2n\sin\frac\pi n$, et le développement
$\sin x = x - \frac{x^3}6 + O(x^5)$ donne
\[
L_n = 2\pi - \frac{\pi^3}{3n^2} + O\Bigl(\frac1{n^4}\Bigr) :
\]
les \omterm{def:b2:curves:length}{longueurs} polygonales du \cref{thm:b2:curves:lengthinv}
convergent \emph{quadratiquement}. Numériquement~: $L_6 = 6$
(l'hexagone, donnant le grossier $\pi > 3$), tandis que $L_{96} =
192\sin\frac{\pi}{96} \approx 6.28206$ contre $2\pi \approx
6.28319$ --- l'erreur $0.00113$ concorde avec la valeur prédite
$\pi^3/(3\cdot96^2) \approx 0.00112$. C'est pourquoi Archimède,
en doublant l'hexagone cinq fois jusqu'à $96$ côtés, pouvait
encadrer $\pi$ à trois décimales à la main~: chaque doublement
divise l'erreur par quatre. La borne supérieure de la
caractérisation polygonale n'est pas seulement atteinte à la
limite~; elle est atteinte \emph{rapidement}, parce qu'une
courbe lisse ne s'écarte de ses cordes qu'au second \omterm{def:b2:structures:generated}{ordre}.
\end{example}

\begin{theorem}[Paramétrage par longueur d'arc]\label{thm:b2:curves:arclength}
Soit $\gamma \colon I \to \R^n$ un arc $\mathcal{C}^k$ régulier
($k \geq 1$). L'\omterm{def:b2:curves:length}{abscisse curviligne} $s$ est un $\mathcal{C}^k$-difféomorphisme
de $I$ sur un intervalle $J$, et $\tilde\gamma =
\gamma \circ s^{-1}$ vérifie $\norm{\tilde\gamma'} = 1$
partout. À translation du paramètre et orientation près,
ce paramétrage \emph{par \omterm{def:b2:curves:length}{longueur d'arc}} (ou \emph{normal}) est
unique.
\end{theorem}

\begin{proof}
$s'(t) = \norm{\gamma'(t)} > 0$ et $s'$ est $\mathcal{C}^{k-1}$
(composée de l'application $\mathcal{C}^{k-1}$ $\gamma'$ avec la
\omterm{def:b2:nvs:norm}{norme}, lisse en dehors de $0$), donc $s$ est $\mathcal{C}^k$,
strictement croissante, une bijection sur $J = s(I)$, et son
inverse est $\mathcal{C}^k$ par le théorème d'inversion locale en
dimension $1$ (volume de première année). Alors
\[
\tilde\gamma'(\sigma)
= \frac{\gamma'(t)}{s'(t)}
= \frac{\gamma'(t)}{\norm{\gamma'(t)}},
\qquad t = s^{-1}(\sigma),
\]
un vecteur \omterm{def:b2:hermitian:adjoint}{unitaire}. Si $\hat\gamma = \gamma\circ\theta$ est un
autre paramétrage normal, alors $\abs{\theta'} = 1$, donc $\theta'
= \pm 1$ constant (continuité), c'est-à-dire $\theta(u) = \pm u + c$.
\end{proof}

\begin{remark}
La \omterm{def:b2:curves:length}{longueur d'arc} est le paramètre qui sépare la géométrie de
la dynamique. Une trajectoire peut être parcourue avec
n'importe quel profil de vitesse --- la physique du mouvement
--- mais toute question invariante par paramétrage (forme,
\omterm{thm:b2:curves:frenet2d}{courbure}, osculation) possède une horloge canonique, la
distance parcourue. C'est pourquoi toutes les formules de
\omterm{thm:b2:curves:frenet2d}{courbure} ci-dessous sont \emph{définies} à vitesse unité puis
\emph{traduites} à des paramétrages quelconques par la règle
de dérivation composée~: les facteurs de traduction sont des
puissances de $v = s'$, et les suivre correctement est tout le
contenu de la \cref{prop:b2:curves:kappaformula}.
\end{remark}

\begin{example}[Cercle et hélice]\label{ex:b2:curves:helix}
Pour le cercle $\gamma(t) = (R\cos t, R\sin t)$, $\norm{\gamma'} =
R$, donc $s = Rt$ et la \omterm{def:b2:curves:length}{longueur} d'un tour \omterm{def:b2:metric:complete}{complet} est $2\pi R$.
Pour l'hélice $\gamma(t) = (a\cos t,\ a\sin t,\ bt)$ avec $a > 0$,
$\norm{\gamma'(t)} = \sqrt{a^2 + b^2}$ est constante~: l'hélice est
parcourue à vitesse constante, et $s = t\sqrt{a^2 + b^2}$.
\end{example}

\begin{example}[Longueur d'arc en coordonnées
polaires]\label{ex:b2:curves:polarlength}
Une courbe polaire $r = r(\theta)$ est l'arc $\gamma(\theta) =
(r\cos\theta,\ r\sin\theta)$, avec
\[
\gamma'(\theta) = (r'\cos\theta - r\sin\theta,\
r'\sin\theta + r\cos\theta),
\qquad
\norm{\gamma'(\theta)}^2 = r'^2 + r^2
\]
(les termes croisés s'annulent)~: l'élément de \omterm{def:b2:curves:length}{longueur} polaire est
$\sqrt{r^2 + r'^2}\,\dd\theta$. Pour la cardioïde $r = 1 +
\cos\theta$~:
\[
r^2 + r'^2 = (1 + \cos\theta)^2 + \sin^2\theta = 2 +
2\cos\theta = 4\cos^2\tfrac\theta2,
\]
et sur $\intcc{-\pi}{\pi}$ le cosinus de l'angle moitié est
positif, donc
\[
L = \int_{-\pi}^{\pi}2\cos\tfrac\theta2\,\dd\theta
= \Bigl[4\sin\tfrac\theta2\Bigr]_{-\pi}^{\pi}
= 4 - (-4) = 8 :
\]
comme l'arche de cycloïde de l'\cref{exo:b2:curves:1}, une courbe
construite à partir de cercles a une \omterm{def:b2:curves:length}{longueur} rationnelle, sans
aucun $\pi$. La factorisation par l'angle moitié est l'astuce
standard pour les \omterm{def:b2:curves:length}{longueurs} de courbes engendrées par des
cercles~; quand elle échoue (l'ellipse), la \omterm{def:b2:curves:length}{longueur} est une
fonction véritablement nouvelle --- une intégrale elliptique,
au-delà des formes closes élémentaires.
\end{example}

\section{Courbure dans le plan}

Tout au long de cette section, les arcs sont $\mathcal{C}^2$ et
réguliers dans le plan euclidien orienté. On paramètre par
\omterm{def:b2:curves:length}{longueur d'arc} et on note $T(s) = \tilde\gamma'(s)$ la \omterm{def:b2:curves:arc}{tangente}
\omterm{def:b2:hermitian:adjoint}{unitaire}, et $N(s)$ le vecteur \omterm{def:b2:hermitian:adjoint}{unitaire} directement orthogonal à
$T(s)$ (rotation de $T$ de $+\pi/2$).

\begin{theorem}[Formules de Frenet dans le plan]\label{thm:b2:curves:frenet2d}
Soit $\tilde\gamma$ un arc $\mathcal{C}^2$ paramétré par \omterm{def:b2:curves:length}{longueur
d'arc} dans le plan orienté. Il existe une fonction \omterm{def:b2:metric:continuity}{continue}
$\kappa$, la \emph{\omterm{thm:b2:curves:frenet2d}{courbure}} (algébrique)\index{courbure}, telle
que
\[
T'(s) = \kappa(s)\, N(s),
\qquad
N'(s) = -\kappa(s)\, T(s).
\]
\end{theorem}

\begin{proof}
Comme $\norm{T(s)}^2 = 1$ pour tout $s$, en dérivant le produit
scalaire on obtient $2\langle T'(s), T(s)\rangle = 0$~: $T'(s)$ est
orthogonal à $T(s)$, donc colinéaire à $N(s)$ (dimension $2$)~; on
écrit $T'(s) = \kappa(s) N(s)$ avec $\kappa(s) = \langle T'(s),
N(s)\rangle$, \omterm{def:b2:metric:continuity}{continue}. De même $N' \perp N$, donc $N' = \lambda
T$~; et en dérivant $\langle T, N\rangle = 0$ on obtient $\langle
T', N\rangle + \langle T, N'\rangle = \kappa + \lambda = 0$.
\end{proof}

\begin{definition}\label{def:b2:curves:curvature}
Lorsque $\kappa(s) \neq 0$, le \emph{rayon de
courbure}\index{rayon de courbure} est $R(s) =
1/\abs{\kappa(s)}$ et le \emph{centre de courbure}\index{centre de
courbure} est
$\tilde\gamma(s) + \frac{1}{\kappa(s)} N(s)$~; le cercle de ce
centre et de rayon $R(s)$ est le \emph{cercle
osculateur}\index{cercle osculateur}, la
meilleure approximation circulaire de la courbe en $\tilde\gamma(s)$.
\end{definition}

\begin{example}[Le cercle osculateur de
l'exponentielle]\label{ex:b2:curves:oscexp}
Pour $y = \eu^x$ au point $(0, 1)$~: $f'(0) = f''(0) = 1$,
donc par la formule pour un graphe ci-dessous,
\[
\kappa(0) = \frac{1}{(1 + 1)^{3/2}} = \frac1{2\sqrt2},
\qquad R = 2\sqrt2 .
\]
La \omterm{def:b2:curves:arc}{tangente} \omterm{def:b2:hermitian:adjoint}{unitaire} est $T = \frac{(1, 1)}{\sqrt2}$, la
normale directe $N = \frac{(-1, 1)}{\sqrt2}$, et le \omterm{def:b2:curves:curvature}{centre de
courbure} est
\[
(0, 1) + 2\sqrt2\cdot\frac{(-1, 1)}{\sqrt2} = (-2,\ 3) :
\]
le \omterm{def:b2:curves:curvature}{cercle osculateur} a pour équation $(x + 2)^2 + (y - 3)^2 =
8$. Pour vérifier l'affirmation de «~meilleure approximation
circulaire~»~: en résolvant l'équation du cercle en $y$ près
de $(0,1)$ et en développant, on obtient $y = 1 + x +
\frac{x^2}2 + O(x^3)$ ---
exactement le développement de Taylor à l'\omterm{def:b2:structures:generated}{ordre} deux de
$\eu^x$. Le \omterm{def:b2:curves:curvature}{cercle osculateur} coïncide en valeur, en pente
\emph{et} en dérivée seconde~; un cercle tangent ordinaire ne
coïnciderait qu'avec les deux premières.
\end{example}

\begin{example}[La développée d'un cercle est son
centre]\label{ex:b2:curves:circleevolute}
Pour le cercle de rayon $R$ parcouru dans le sens direct,
$\kappa = 1/R$ et $N$ pointe vers le centre, donc le
\omterm{def:b2:curves:curvature}{centre de courbure} $\tilde\gamma + \frac1\kappa N$ est le
centre du cercle, pour tout $s$~: le \omterm{def:b2:curves:curvature}{cercle osculateur}
d'un cercle est le cercle lui-même, et le lieu des \omterm{def:b2:curves:curvature}{centres
de courbure} se réduit à un point. Ce cas dégénéré
calibre l'\cref{exo:b2:curves:6}~: là, la vitesse de la
développée est $-\frac{\kappa'}{\kappa^2}N$, qui s'annule
identiquement précisément lorsque $\kappa$ est constante.
\end{example}

\begin{proposition}[Courbure dans un paramétrage
quelconque]\label{prop:b2:curves:kappaformula}
Pour un arc plan $\mathcal{C}^2$ régulier $\gamma(t) = (x(t), y(t))$,
\[
\kappa(t)
= \frac{x'(t)\,y''(t) - y'(t)\,x''(t)}
       {\bigl(x'(t)^2 + y'(t)^2\bigr)^{3/2}} ,
\]
en particulier $\kappa = \dfrac{y''}{(1 + y'^2)^{3/2}}$ pour un
graphe $y = f(x)$.
\end{proposition}

\begin{proof}
Posons $v(t) = \norm{\gamma'(t)} = s'(t)$, de sorte que $\gamma' = vT$
(en composant les données à vitesse unité avec $s$). En dérivant,
\[
\gamma'' = v'T + v\,T'\cdot s' = v'T + v^2\kappa N .
\]
Prenons maintenant le \omterm{def:b2:linalg:det}{déterminant} (dans la base orientée
canonique) de
$(\gamma', \gamma'')$~: comme $\det(T, T) = 0$ et $\det(T, N) = 1$,
\[
\det(\gamma', \gamma'') = \det(vT,\ v'T + v^2\kappa N)
= v^3\kappa .
\]
Le membre de gauche est $x'y'' - y'x''$, et $v^3 = (x'^2 +
y'^2)^{3/2}$. Le cas du graphe est le paramétrage $t \mapsto (t,
f(t))$.
\end{proof}

\begin{example}[Cercle, droite, parabole]\label{ex:b2:curves:kappaexamples}
Une droite a $\kappa = 0$ (et réciproquement~: $T' = 0$ signifie
$T$ constant, donc $\tilde\gamma(s) = \tilde\gamma(0) + sT$, une
droite). Le cercle de rayon $R$ parcouru dans le sens direct a
$\kappa = 1/R$~: avec $\gamma(t) = (R\cos t, R\sin t)$, la formule
donne $\kappa = R^2/R^3$. Pour la parabole $y = x^2/2$~:
$\kappa(x) = 1/(1 + x^2)^{3/2}$, maximale au sommet --- la parabole
est le plus fortement courbée là où elle fait demi-tour.
\end{example}

\begin{remark}[Pièges courants autour de la courbure]
(i) La \omterm{thm:b2:curves:frenet2d}{courbure} algébrique d'un arc plan change de signe lorsqu'on
inverse l'orientation de l'arc ou du plan~: seuls
$\abs\kappa$ et $R = 1/\abs\kappa$ sont purement géométriques. Un
cercle parcouru dans le sens rétrograde a $\kappa = -1/R$. (ii) La
formule pour un graphe $\kappa = f''/(1 + f'^2)^{3/2}$ choisit
tacitement le paramétrage par $x$~; l'appliquer à une courbe
qui n'est pas un graphe près du point (\omterm{def:b2:curves:arc}{tangente} verticale) est la
bévue classique. (iii) En un point où $\gamma' = 0$,
rien n'est défini --- ni $T$ ni $\kappa$ --- et la
trajectoire peut véritablement se briser
(\cref{ex:b2:curves:cusp})~; toujours vérifier la régularité avant
de dériver la \omterm{def:b2:curves:arc}{tangente} \omterm{def:b2:hermitian:adjoint}{unitaire}. (iv) Dans l'espace, $\kappa =
\norm{T'} \geq 0$ par convention~: il n'y a pas de signe à se
tromper, mais pas non plus de signe à exploiter --- l'information
de type inflexion passe dans la \omterm{thm:b2:curves:frenet3d}{torsion}. (v) Enfin, $\kappa$ est
une dérivée \emph{par rapport à la \omterm{def:b2:curves:length}{longueur d'arc}}~: pour un
paramétrage qui n'est pas à vitesse unité, oublier le facteur $v^3$
dans la \cref{prop:b2:curves:kappaformula} est l'erreur la plus
fréquente en pratique.
\end{remark}

\begin{figure}[ht]
\centering
\begin{tikzpicture}[scale=1.6]
  % parabola y = x^2/2
  \draw[->] (-2.2,0) -- (2.4,0) node[below] {$x$};
  \draw[->] (0,-0.4) -- (0,2.4) node[left] {$y$};
  \draw[thick, ocDarkBlue, domain=-2.1:2.1, smooth, samples=60]
    plot (\x, {0.5*\x*\x});
  % osculating circle at vertex: kappa = 1 => R = 1, center (0,1)
  \draw[ocRed, dashed] (0,1) circle (1);
  \fill[ocRed] (0,1) circle (0.03) node[right, font=\small]
    {centre de courbure};
  \fill (0,0) circle (0.03);
  % tangent and normal at vertex
  \draw[->, thick] (0,0) -- (0.7,0) node[below, font=\small] {$T$};
  \draw[->, thick] (0,0) -- (0,0.7) node[left, font=\small] {$N$};
  % a second point
  \fill (1.2,0.72) circle (0.03);
  \draw[->, thick] (1.2,0.72) -- ++(0.448,0.5376)
    node[right, font=\small] {$T$};
  \draw[->, thick] (1.2,0.72) -- ++(-0.5376,0.448)
    node[above, font=\small] {$N$};
\end{tikzpicture}
\caption{La parabole $y = x^2/2$, son repère de Frenet mobile $(T,
N)$, et le \omterm{def:b2:curves:curvature}{cercle osculateur} au sommet (rayon $1$, puisque
$\kappa(0) = 1$). Le repère tourne à mesure que le point se
déplace~; la \omterm{thm:b2:curves:frenet2d}{courbure} est le taux de cette rotation par unité de
\omterm{def:b2:curves:length}{longueur d'arc}.}
\label{fig:b2:curves:osculating}
\end{figure}

\begin{theorem}[La courbure détermine la courbe]\label{thm:b2:curves:fundamental}
Soit $\kappa \colon J \to \R$ \omterm{def:b2:metric:continuity}{continue}. Il existe un arc
$\mathcal{C}^2$ du plan paramétré par \omterm{def:b2:curves:length}{longueur d'arc} de \omterm{thm:b2:curves:frenet2d}{courbure}
$\kappa$, et il est unique à une isométrie directe près (rotation
suivie d'une translation).
\end{theorem}

\begin{proof}
\emph{Existence.} Fixons $s_0 \in J$ et posons $\varphi(s) =
\int_{s_0}^s \kappa(u)\,\dd u$, puis
\[
\tilde\gamma(s) = \Bigl(\int_{s_0}^s \cos\varphi(u)\,\dd u,\
\int_{s_0}^s \sin\varphi(u)\,\dd u\Bigr).
\]
Alors $T(s) = (\cos\varphi(s), \sin\varphi(s))$ est un vecteur
\omterm{def:b2:hermitian:adjoint}{unitaire}, $N(s) = (-\sin\varphi, \cos\varphi)$, et
\[
T'(s) = \varphi'(s)\,(-\sin\varphi, \cos\varphi)
= \kappa(s)\,N(s) :
\]
l'arc est paramétré par \omterm{def:b2:curves:length}{longueur d'arc}, de \omterm{thm:b2:curves:frenet2d}{courbure} $\kappa$.

\emph{Unicité.} Soient $\gamma_1, \gamma_2$ deux \omterm{def:b2:curves:arc}{arcs paramétrés}
par \omterm{def:b2:curves:length}{longueur d'arc} de même \omterm{thm:b2:curves:frenet2d}{courbure}. Chaque \omterm{def:b2:curves:arc}{tangente} \omterm{def:b2:hermitian:adjoint}{unitaire} se
relève en une fonction angle,
par une construction explicite~: on voit $T_j$ comme le
nombre complexe $z_j = a_j + \iu b_j$ de module $1$, on choisit
$\varphi_j(0)$ tel que $z_j(0) = \eu^{\iu\varphi_j(0)}$, et l'on pose
\[
\varphi_j(s) = \varphi_j(0) + \int_0^s\det\bigl(T_j,
T_j'\bigr)(u)\,\dd u .
\]
De $\abs{z_j} = 1$~: $\operatorname{Re}(\conj{z_j}\,z_j') =
0$, donc $\conj{z_j}\,z_j' = \iu\det(T_j, T_j') =
\iu\,\varphi_j'$, c'est-à-dire $z_j' = \iu\varphi_j'z_j$~; alors
\[
\bigl(z_j\,\eu^{-\iu\varphi_j}\bigr)'
= \eu^{-\iu\varphi_j}\bigl(z_j' - \iu\varphi_j'z_j\bigr) = 0,
\]
donc $z_j = \eu^{\iu\varphi_j}$ partout~: $T_j =
(\cos\varphi_j, \sin\varphi_j)$ avec $\varphi_j$ de classe
$\mathcal C^1$. De plus $\det(T_j, T_j') = \det(T_j,
\kappa N_j) = \kappa$, donc $\varphi_j' = \kappa$. D'où $\varphi_2 = \varphi_1 + c$ pour une
constante $c$~: $T_2$ est $T_1$ tourné de l'angle fixe $c$, donc
en intégrant, $\gamma_2 = \rho(\gamma_1) + w$ où $\rho$ est la
rotation d'angle $c$ et $w$ un vecteur constant.
\end{proof}

\begin{remark}
C'est le prototype unidimensionnel d'un \emph{théorème fondamental
de la géométrie}~: un ensemble \omterm{def:b2:metric:complete}{complet} d'invariants locaux (ici,
une fonction) classifie l'objet à un déplacement près. La
version tridimensionnelle ci-dessous nécessite deux invariants.
\end{remark}

\begin{example}[Courbure constante signifie
cercle]\label{ex:b2:curves:constantkappa}
Prenons $\kappa \equiv \kappa_0 > 0$ dans la formule d'existence~:
$\varphi(s) = \kappa_0 s$ et
\[
\tilde\gamma(s) = \Bigl(\frac{\sin\kappa_0s}{\kappa_0},\
\frac{1 - \cos\kappa_0s}{\kappa_0}\Bigr) :
\]
le cercle de rayon $1/\kappa_0$ centré en $(0,
1/\kappa_0)$, parcouru à vitesse unité. Par la moitié
«~unicité~» du théorème, \emph{tout} \omterm{def:b2:curves:arc}{arc paramétré} par \omterm{def:b2:curves:length}{longueur
d'arc} de \omterm{thm:b2:curves:frenet2d}{courbure} constante $\kappa_0$ est un morceau d'un cercle
de rayon $1/\kappa_0$ (ou une droite si $\kappa_0 = 0$) --- la
réciproque du calcul de l'\cref{ex:b2:curves:kappaexamples}, et
le cas plan de l'\cref{exo:b2:curves:9}.
\end{example}

\begin{example}[Reconstruire une courbe à partir de sa
courbure]\label{ex:b2:curves:reconstruction}
Quelle courbe paramétrée par \omterm{def:b2:curves:length}{longueur d'arc} a pour \omterm{def:b2:curves:curvature}{rayon de
courbure} $R(s) = 1 +
s^2$~? En suivant la preuve d'existence avec $\kappa(s) =
\frac1{1+s^2}$ et $s_0 = 0$~: $\varphi(s) = \arctan s$, donc
\[
T(s) = (\cos\arctan s,\ \sin\arctan s)
= \Bigl(\frac{1}{\sqrt{1+s^2}},\
\frac{s}{\sqrt{1+s^2}}\Bigr),
\]
et en intégrant,
\[
\tilde\gamma(s) = \Bigl(\ln\bigl(s + \sqrt{1 + s^2}\bigr),\
\sqrt{1 + s^2} - 1\Bigr).
\]
En posant $x = \ln(s + \sqrt{1+s^2})$, c'est-à-dire $s = \sinh x$, la
seconde coordonnée est $\cosh x - 1$~: la courbe est la
\emph{chaînette} $y = \cosh x - 1$. Cela boucle la boucle avec
l'\cref{exo:b2:curves:3}, où nous avons calculé $R = \cosh^2 x = 1 +
\sinh^2 x = 1 + s^2$ directement~: le théorème fondamental
garantit que la chaînette est la \emph{seule} courbe ayant ce
profil de \omterm{thm:b2:curves:frenet2d}{courbure}, à une isométrie directe près.
\end{example}

\begin{remark}[Où la courbure sert ensuite]
La décomposition $\gamma'' = v'T + v^2\kappa N$ obtenue dans
la preuve de la \cref{prop:b2:curves:kappaformula} est la
cinématique de tout mouvement courbe~: accélération tangentielle
contre accélération centripète. La \omterm{thm:b2:curves:frenet2d}{courbure} revient pour les
surfaces
(\cref{ch:b2:surfaces}) à travers la \omterm{thm:b2:curves:frenet2d}{courbure} des courbes qu'on
y trace, et le calcul des \omterm{pb:b2:curves:1}{enveloppes} du problème de week-end de
ce chapitre --- développées, caustiques --- est l'optique
géométrique des fronts d'onde. Le volume de troisième année
reprend le point de vue intrinsèque pour les sous-variétés de
$\R^n$.
\end{remark}

\section{Repère de Frenet dans l'espace}

Soit maintenant $\tilde\gamma \colon J \to \R^3$ un arc
$\mathcal{C}^3$ paramétré par \omterm{def:b2:curves:length}{longueur d'arc} qui est
\emph{birégulier}~: $T'(s) \neq 0$ pour
tout $s$. Alors $\kappa(s) = \norm{T'(s)} > 0$ définit la
\emph{\omterm{thm:b2:curves:frenet2d}{courbure}} (pas de signe dans l'espace~: il n'y a pas
d'orientation privilégiée de la normale), et l'on pose~:
\[
N(s) = \frac{T'(s)}{\kappa(s)}
\quad\text{(normale principale)},
\qquad
B(s) = T(s) \wedge N(s)
\quad\text{(binormale)} ,
\]
de sorte que $(T, N, B)$ est un repère orthonormé direct, le
\emph{repère de Frenet}. Le plan passant par $\tilde\gamma(s)$
engendré par $T, N$ est le \emph{plan osculateur}.

\begin{theorem}[Formules de Frenet dans l'espace]\label{thm:b2:curves:frenet3d}
Il existe une fonction \omterm{def:b2:metric:continuity}{continue} $\tau$, la \emph{torsion}\index{torsion}, telle que
\[
T' = \kappa N, \qquad
N' = -\kappa T + \tau B, \qquad
B' = -\tau N .
\]
\end{theorem}

\begin{proof}
La première formule est la définition de $N$. Chacun des vecteurs
$T, N, B$ est de \omterm{def:b2:nvs:norm}{norme} constante $1$ et ils sont deux à deux
orthogonaux~; en dérivant les six relations $\langle X, Y\rangle =
\delta_{XY}$, on voit que la matrice de $(T', N', B')$ dans la base
$(T, N, B)$ est antisymétrique~: en effet $\langle X', Y\rangle +
\langle X, Y'\rangle = 0$ et
$\langle X', X\rangle = 0$. Son coefficient $(N, T)$ est $\langle
N', T\rangle = -\langle N, T'\rangle = -\kappa$, et son coefficient
de colonne $(T, B)$ est
$\langle T', B\rangle = \kappa\langle N, B\rangle = 0$. En nommant
$\tau = \langle N', B\rangle$ le coefficient libre restant, on
obtient exactement les trois formules affichées~: l'antisymétrie
complète $\langle B', N\rangle = -\tau$ et
$\langle B', T\rangle = 0$. La continuité de $\tau = \langle N',
B\rangle$ est claire puisque $N'$ et $B$ sont \omterm{def:b2:metric:continuity}{continus}
($\tilde\gamma$ est
$\mathcal{C}^3$, donc $N = T'/\kappa$ est $\mathcal{C}^1$).
\end{proof}

\begin{example}[Le vecteur de Darboux]\label{ex:b2:curves:darboux}
Les trois formules de Frenet se condensent en une seule. Posons
$\omega(s) =
\tau\,T + \kappa\,B$ (le \emph{vecteur de Darboux}). En utilisant
$B \wedge T = N$, $T \wedge N = B$, $N \wedge B = T$~:
\[
\omega \wedge T = \kappa\,N = T', \qquad
\omega \wedge N = \tau\,B - \kappa\,T = N', \qquad
\omega \wedge B = -\tau\,N = B' :
\]
chaque vecteur du repère évolue selon $X' = \omega \wedge X$, la
signature cinématique d'une \emph{rotation instantanée} de
vecteur vitesse angulaire $\omega$. Le repère tourne au taux
$\norm\omega = \sqrt{\kappa^2 + \tau^2}$ autour de l'axe mobile
$\omega$~; la \omterm{thm:b2:curves:frenet2d}{courbure} est la composante de la rotation autour de
la binormale, la \omterm{thm:b2:curves:frenet3d}{torsion} la composante autour de la \omterm{def:b2:curves:arc}{tangente}. Pour
l'hélice, $\omega$ est un vecteur constant le long de l'axe du
cylindre --- ce qui est exactement pourquoi le repère de l'hélice
précesse régulièrement. L'antisymétrie de la matrice de Frenet,
exploitée dans l'\cref{exo:b2:curves:7}, est la forme matricielle
de ce seul fait géométrique.
\end{example}

\begin{proposition}[La torsion mesure la planéité]\label{prop:b2:curves:torsion}
Un arc birégulier est contenu dans un plan si et seulement si
$\tau \equiv 0$~; dans ce cas le plan est le plan osculateur
(constant).
\end{proposition}

\begin{proof}
Si $\tau \equiv 0$, alors $B' = 0$, donc $B$ est un vecteur
\omterm{def:b2:hermitian:adjoint}{unitaire} constant $B_0$, et
\[
\frac{\dd}{\dd s}\langle \tilde\gamma(s), B_0\rangle
= \langle T(s), B_0\rangle = 0 :
\]
$\langle \tilde\gamma, B_0\rangle$ est constant, donc l'arc est
contenu dans un plan orthogonal à $B_0$. Réciproquement, si l'arc
est contenu dans un plan $P$,
alors $T$ et $T'$ (donc $N$) sont parallèles à la direction de
$P$ pour tout $s$~; donc $B = T \wedge N$ est l'une des deux
normales \omterm{def:b2:hermitian:adjoint}{unitaires} de $P$, et étant \omterm{def:b2:metric:continuity}{continue} elle est constante~;
alors $0 = B' = -\tau N$ avec $N \neq 0$ force $\tau \equiv 0$.
\end{proof}

\begin{example}[Un cercle incliné a une torsion nulle]
L'arc $\gamma(t) = \bigl(\cos t,\ \tfrac{\sin t}{\sqrt2},\
\tfrac{\sin t}{\sqrt2}\bigr)$ est contenu dans le plan $y = z$, et
est le cercle unité de ce plan (vérifier~: $\norm{\gamma(t)} =
1$ et la base orthonormée $(1,0,0)$,
$(0, \tfrac1{\sqrt2}, \tfrac1{\sqrt2})$ du plan met en évidence le
paramétrage standard). Sans aucun calcul de Frenet,
la \cref{prop:b2:curves:torsion} prédit $\tau \equiv 0$, et
la binormale fixe doit être la normale \omterm{def:b2:hermitian:adjoint}{unitaire} $\pm(0,
\tfrac1{\sqrt2}, -\tfrac1{\sqrt2})$ du plan. La \omterm{thm:b2:curves:frenet3d}{torsion} ne mesure
pas le fait d'être «~incliné dans l'espace~»~; elle mesure le fait
de \emph{quitter} un plan. Seul le $b$ non nul de l'hélice
ci-dessous produit une véritable \omterm{thm:b2:curves:frenet3d}{torsion}.
\end{example}

\begin{example}[L'hélice]\label{ex:b2:curves:helixfrenet}
Pour l'hélice $\gamma(t) = (a\cos t, a\sin t, bt)$, $a > 0$, nous
avons calculé $s = ct$ avec $c = \sqrt{a^2 + b^2}$. Alors
\[
T = \frac1c(-a\sin t,\ a\cos t,\ b),
\qquad
T' \cdot \frac{\dd t}{\dd s}
= \frac{1}{c^2}(-a\cos t, -a\sin t, 0),
\]
donc $\kappa = a/c^2 = a/(a^2 + b^2)$ et $N = (-\cos t, -\sin t,
0)$~: la normale principale pointe horizontalement vers l'axe.
Ensuite $B = T \wedge N = \frac1c(b\sin t, -b\cos t, a)$, et $B'
\frac{\dd t}{\dd s} = \frac{b}{c^2}(\cos t, \sin t, 0) = -\tau N$
donne
\[
\boxed{\ \kappa = \frac{a}{a^2 + b^2}, \qquad
\tau = \frac{b}{a^2 + b^2}. \ }
\]
Les deux invariants sont constants --- et l'on peut montrer,
réciproquement, que les seules courbes birégulières à $\kappa > 0$
constant et $\tau$ constant sont les hélices (des cercles lorsque
$\tau = 0$). Notons les signes~: $b >
0$ donne une hélice dextre à \omterm{thm:b2:curves:frenet3d}{torsion} positive.
\end{example}

\begin{example}[Courbure et torsion sans longueur d'arc~; la
cubique gauche]\label{ex:b2:curves:twistedcubic}
Reparamétrer par \omterm{def:b2:curves:length}{longueur d'arc} est en général impossible sous
forme close, si bien que les invariants doivent être extraits des
dérivées brutes. Posons $v = \norm{\gamma'} = s'$~; alors $\gamma' =
vT$ et, comme dans la preuve de la
\cref{prop:b2:curves:kappaformula},
\[
\gamma'' = v'T + v^2\kappa N,
\qquad
\gamma' \wedge \gamma'' = v^3\kappa\,(T \wedge N) =
v^3\kappa\,B .
\]
En prenant les \omterm{def:b2:nvs:norm}{normes} ($\kappa \geq 0$ dans l'espace)~:
\[
\kappa = \frac{\norm{\gamma' \wedge \gamma''}}{v^3} .
\]
En dérivant $\gamma''$ une fois de plus et en convertissant $N' =
v(-\kappa T + \tau B)$ (règle de dérivation composée à travers
$s$), la seule
composante en $B$ provient du dernier terme~:
\[
\gamma''' = \bigl(v'' - v^3\kappa^2\bigr)T +
\bigl(v'v\kappa + (v^2\kappa)'\bigr)N + v^3\kappa\tau\,B,
\]
de sorte que, en appariant avec $\gamma' \wedge \gamma'' =
v^3\kappa B$,
\[
\det(\gamma', \gamma'', \gamma''') = \langle\gamma' \wedge
\gamma'',\ \gamma'''\rangle = v^6\kappa^2\tau,
\qquad\text{c'est-à-dire}\qquad
\tau = \frac{\det(\gamma', \gamma'',
\gamma''')}{\norm{\gamma' \wedge \gamma''}^2} .
\]
Application à la \emph{cubique gauche} $\gamma(t) = (t,\
t^2,\ t^3)$ en $t = 0$~: $\gamma' = (1, 0, 0)$, $\gamma'' =
(0, 2, 0)$, $\gamma''' = (0, 0, 6)$, donc $v = 1$,
\[
\gamma' \wedge \gamma'' = (0, 0, 2), \qquad \kappa(0) = 2,
\qquad
\det(\gamma', \gamma'', \gamma''') = 12, \qquad \tau(0) =
\frac{12}{4} = 3 .
\]
Éclairage final~: les deux formules sont des rapports où la vitesse
$v$ se simplifie exactement au degré nécessaire --- $\kappa$ se
comporte comme une dérivée seconde par unité de \omterm{def:b2:curves:length}{longueur}, $\tau$
comme le volume mixte de trois dérivées par aire au carré --- ce
qui explique \emph{pourquoi} elles sont géométriques alors que
$\gamma''$ lui-même ne l'est pas.
\end{example}

\begin{remark}[Théorème fondamental pour les courbes de l'espace]
Comme dans le plan, le couple $(\kappa, \tau)$ avec $\kappa > 0$
détermine un arc birégulier à une isométrie directe de $\R^3$
près~: les formules de Frenet forment un système différentiel
linéaire pour le repère
$(T, N, B)$, auquel s'applique la théorie de Cauchy--Lipschitz du
\cref{ch:b2:diffeq}~; l'orthonormalité du repère solution est
préservée parce que la matrice des coefficients est antisymétrique
(le même argument de matrice de Gram que dans
l'\cref{exo:b2:curves:7}), et la courbe
est retrouvée en intégrant $T$. Nous laissons les détails au
lecteur comme un exercice substantiel mais instructif.
\end{remark}

\section{Étude locale~: position par rapport à la tangente}

\begin{proposition}[Allure locale en un point
régulier]\label{prop:b2:curves:local}
Soit $\gamma$ un arc plan de classe $\mathcal{C}^k$ en $t_0$, avec
$p$ le plus petit indice tel que $\gamma^{(p)}(t_0) \neq 0$ et $q$ le
plus petit indice $> p$ tel que $\gamma^{(q)}(t_0)$ ne soit pas
colinéaire à
$\gamma^{(p)}(t_0)$ (en supposant que les deux existent, $q \leq
k$). Dans la base
$(u, v) = (\gamma^{(p)}(t_0), \gamma^{(q)}(t_0))$ centrée en
$\gamma(t_0)$, la formule de Taylor--Young donne les coordonnées
\[
X(t) \sim \frac{(t - t_0)^p}{p!},
\qquad
Y(t) \sim \frac{(t - t_0)^q}{q!} .
\]
L'allure locale ne dépend que des parités de $p$ et $q$~:
\begin{center}
\begin{tabular}{c p{3.5cm} p{5.7cm}}
$p$ impair, $q$ pair & \emph{point ordinaire} & la courbe ne
traverse aucune droite, reste d'un côté de la \omterm{def:b2:curves:arc}{tangente}\\
$p$ impair, $q$ impair & \emph{point d'inflexion} & la courbe
traverse sa \omterm{def:b2:curves:arc}{tangente}\\
$p$ pair, $q$ impair & \emph{point de rebroussement de première
espèce} & les deux branches de part et d'autre de la \omterm{def:b2:curves:arc}{tangente}\\
$p$ pair, $q$ pair & \emph{point de rebroussement de seconde
espèce} & les deux branches du même côté\\
\end{tabular}
\end{center}
\end{proposition}

\begin{proof}
Taylor--Young à l'\omterm{def:b2:structures:generated}{ordre} $q$ (la fonction $\gamma$ est
$\mathcal{C}^q$ près de $t_0$)~:
\[
\gamma(t) - \gamma(t_0)
= \sum_{j=p}^{q} \frac{(t-t_0)^j}{j!}\,\gamma^{(j)}(t_0)
+ o\bigl((t-t_0)^q\bigr).
\]
Par le choix de $p$ et $q$, chaque $\gamma^{(j)}(t_0)$ avec $p
\leq j < q$ est colinéaire à $u$~; en regroupant les composantes
dans la base $(u, v)$~: $X(t) = \frac{(t-t_0)^p}{p!}(1 + o(1))$ et
$Y(t) = \frac{(t-t_0)^q}{q!}(1 + o(1))$. Le tableau des signes de
$X$ et $Y$ pour $t \gtrless t_0$ --- gouverné exactement par les
parités --- donne les quatre allures~: par exemple si $p$ est pair,
$X > 0$ des deux côtés (les deux branches partent dans la direction
$+u$~: un rebroussement),
et le côté de la \omterm{def:b2:curves:arc}{tangente} ($\operatorname{sign} Y$) bascule
avec $q$ impair.
\end{proof}

\begin{example}
Pour $\gamma(t) = (t^2, t^3)$ en $t_0 = 0$
(\cref{ex:b2:curves:cusp})~: $\gamma'' (0)= (2, 0)$, $\gamma'''(0) =
(0, 6)$, donc $p = 2$, $q = 3$~: un point de rebroussement de
première espèce, l'allure familière de la parabole semi-cubique.
Pour $\gamma(t) = (t,
t^3)$ en $0$~: $p = 1$, $q = 3$~: inflexion --- la cubique traverse
sa \omterm{def:b2:curves:arc}{tangente}.
\end{example}

\begin{remark}[Perspectives au sein de ce volume]
Les courbes alimentent les chapitres suivants de trois façons.
Tracées sur une
surface, elles définissent ses plans tangents et sa première
forme fondamentale (\cref{ch:b2:surfaces}), et leurs \omterm{def:b2:curves:length}{longueurs}
se calculent en restreignant la métrique ambiante --- le
chapitre à venir est en grande partie ce chapitre relativisé. Le
calcul des \omterm{pb:b2:curves:1}{enveloppes} du problème de week-end rencontre les
intégrales doubles au \cref{ch:b2:multint}, où l'aire de
l'\omterm{pb:b2:curves:1}{astroïde} est recalculée par la formule de Green
(\cref{exo:b2:multint:5}) --- une courbe, deux théories, des
réponses concordantes. Et le système de Frenet a déjà utilisé les
équations \omterm{def:b2:diffcalc:differential}{différentielles} linéaires du \cref{ch:b2:diffeq}
(existence, unicité, et l'argument de préservation de
l'orthogonalité de l'\cref{exo:b2:curves:7})~: le théorème
fondamental des courbes est un théorème d'équations
\omterm{def:b2:diffcalc:differential}{différentielles} habillé en géométrie.
\end{remark}

\begin{remark}[Méthode~: mener l'étude locale]
En pratique la classification est une routine en quatre étapes.
\emph{Un}, dériver en $t_0$ jusqu'à ce qu'apparaisse la première
dérivée non nulle~: son indice est $p$, sa valeur le vecteur
$u$. \emph{Deux}, continuer à dériver jusqu'à ce qu'apparaisse une
dérivée non colinéaire à $u$~: indice $q$, vecteur $v$.
\emph{Trois}, lire les parités $(p, q)$ dans le tableau.
\emph{Quatre}, tracer~: la courbe part selon $+u$ si $p$ est impair
(selon $u$ puis en revenant selon $u$ si $p$ est pair), du côté
de $v$ dicté par le signe de $Y$. Deux précautions. Le repère
$(u, v)$ n'est en général \emph{pas} orthonormé --- le tableau
décrit des positions relatives à la \omterm{def:b2:curves:arc}{tangente}, non des angles
ou des distances, donc ne pas lire la \omterm{thm:b2:curves:frenet2d}{courbure} sur l'image.
Et des dérivées intermédiaires colinéaires à $u$ sont permises
entre les rangs $p$ et $q$ (elles ne font que décaler le
développement de $X$)~; ce qui ne doit \emph{pas} arriver, c'est de
s'arrêter à la première dérivée non nulle et de deviner $q = p + 1$~:
pour $\gamma(t) = (t^2, t^4 + t^5)$ la supposition naïve $q = 3$
est fausse, car $\gamma^{(3)}(0)$ est encore colinéaire à
$\gamma''(0)$ --- c'est précisément l'\cref{exo:b2:curves:5}.
\end{remark}

\section{Exercices}

\begin{exercise}[$\star$]\label{exo:b2:curves:1}
Calculer la \omterm{def:b2:curves:length}{longueur} d'une arche de la cycloïde $\gamma(t) = (t -
\sin t,\ 1 - \cos t)$, $t \in [0, 2\pi]$. \emph{(Utiliser $1 - \cos t =
2\sin^2(t/2)$.)}
\end{exercise}

\begin{exercise}[$\star$]\label{exo:b2:curves:2}
Calculer la \omterm{thm:b2:curves:frenet2d}{courbure} de l'ellipse $\gamma(t) = (a\cos t,\ b\sin
t)$ ($a > b > 0$) et localiser les points de \omterm{thm:b2:curves:frenet2d}{courbure} maximale et
minimale.
\end{exercise}

\begin{exercise}[$\star$]\label{exo:b2:curves:3}
Montrer que la \omterm{def:b2:curves:length}{longueur d'arc} du graphe de $f(x) = \cosh x$ sur
$[0, x]$ vaut $\sinh x$, et calculer la \omterm{thm:b2:curves:frenet2d}{courbure} de cette
courbe (la \emph{chaînette}). Vérifier que $R(x) = 1/\kappa(x) =
\cosh^2 x$.
\end{exercise}

\begin{exercise}[$\star\star$]\label{exo:b2:curves:4}
(Spirale logarithmique) Soit $\gamma(t) = e^{t}(\cos t,\ \sin t)$, $t
\in \R$.
Montrer que l'angle entre $\gamma(t)$ et $\gamma'(t)$ est
constant, calculer la \omterm{def:b2:curves:length}{longueur d'arc} de $\gamma$ sur $(-\infty, 0]$
(finie~!), et la \omterm{thm:b2:curves:frenet2d}{courbure}.
\end{exercise}

\begin{exercise}[$\star\star$]\label{exo:b2:curves:5}
Déterminer $p$, $q$ et l'allure locale (point ordinaire,
inflexion, rebroussement) de $\gamma(t) = (t^2,\ t^4 + t^5)$ en $t = 0$, et de
$\gamma(t) = (t^3,\ t^4)$ en $t = 0$.
\end{exercise}

\begin{exercise}[$\star\star$]\label{exo:b2:curves:6}
Soit $\gamma$ un arc plan paramétré par \omterm{def:b2:curves:length}{longueur d'arc} avec
$\kappa(s) > 0$ pour
tout $s$, et soit $c(s) = \gamma(s) + \frac{1}{\kappa(s)}N(s)$ le
\omterm{def:b2:curves:curvature}{centre de courbure} (la courbe $c$ est la \emph{développée}).
En supposant $\kappa$ de classe $\mathcal{C}^1$, montrer que $c'(s) =
-\frac{\kappa'(s)}{\kappa(s)^2}N(s)$~: la développée est \omterm{def:b2:curves:arc}{tangente}
aux normales de $\gamma$.
\end{exercise}

\begin{exercise}[$\star\star\star$]\label{exo:b2:curves:7}
Soit $A(s)$ une famille \omterm{def:b2:metric:continuity}{continue} de matrices antisymétriques $3
\times 3$ et $F' = F A$ une solution matricielle avec $F(s_0)$
orthogonale. Montrer que $F(s)$ est orthogonale pour tout $s$.
\emph{(Dériver $G = F F^{\mathsf T}$ et utiliser l'unicité dans
Cauchy--Lipschitz.)} Expliquer la pertinence pour le système de
Frenet.
\end{exercise}

\begin{exercise}[$\star\star\star$]\label{exo:b2:curves:8}
(\omterm{thm:b2:curves:frenet2d}{Courbure} totale d'une courbe fermée \omterm{def:b2:affine:convex}{convexe}) Soit $\tilde\gamma$ un
arc plan fermé $\mathcal{C}^2$ paramétré par \omterm{def:b2:curves:length}{longueur d'arc}, de
\omterm{def:b2:curves:length}{longueur} $L$
(de sorte que $\tilde\gamma(s + L) = \tilde\gamma(s)$), parcouru
une fois dans le sens direct. En utilisant la fonction angle
$\varphi$ avec $T =
(\cos\varphi, \sin\varphi)$ du
\cref{thm:b2:curves:fundamental}, expliquer pourquoi $\varphi(L) -
\varphi(0)$ est un multiple de $2\pi$, et montrer que
$\int_0^L \kappa(s)\,\dd s = \varphi(L) - \varphi(0)$. (Pour un
cercle de rayon $R$~: $\int \kappa = \frac1R \cdot 2\pi R = 2\pi$.
Le théorème des \omterm{def:b2:curves:arc}{tangentes} tournantes affirme la valeur $2\pi$ pour
toute courbe fermée simple~; on ne vous demande pas de le prouver.)
\end{exercise}

\begin{exercise}[$\star\star\star$]\label{exo:b2:curves:9}
Montrer qu'une courbe gauche birégulière à $\kappa > 0$ constant et
$\tau = 0$ est (un arc d')un cercle de rayon $1/\kappa$.
\emph{(Utiliser la \cref{prop:b2:curves:torsion}, puis montrer que
le centre $\gamma + \frac1\kappa N$ est constant.)}
\end{exercise}

\begin{exercise}[$\star$]\label{exo:b2:curves:10}
Calculer la \omterm{def:b2:curves:length}{longueur d'arc} de la parabole $y = x^2/2$ sur
$\intcc0a$ et montrer qu'elle vaut
\[
\tfrac12\Bigl(a\sqrt{1 + a^2} + \ln\bigl(a + \sqrt{1 +
a^2}\bigr)\Bigr).
\]
\end{exercise}

\begin{exercise}[$\star\star$]\label{exo:b2:curves:11}
Soit $\gamma$ un arc $\mathcal C^2$ régulier de $\R^n$ dont toutes
les \omterm{def:b2:curves:arc}{tangentes} passent par un point fixe $P$. Prouver que
la trajectoire de $\gamma$ est contenue dans une droite.
\emph{(Paramétrer par \omterm{def:b2:curves:length}{longueur d'arc}, écrire $\gamma(s) +
\lambda(s)T(s) = P$ et dériver.)}
\end{exercise}

\begin{exercise}[$\star\star\star$]\label{exo:b2:curves:12}
(Théorème fondamental pour les courbes de l'espace) Soient $\kappa
> 0$ et $\tau$ des fonctions \omterm{def:b2:metric:continuity}{continues} sur un intervalle $J$. Mener
à bien le programme de la remarque qui suit
l'\cref{ex:b2:curves:helixfrenet}~: (a) montrer que le système
linéaire $F' = FA(s)$, avec $A(s)$ la matrice de Frenet
antisymétrique construite à partir de $\kappa, \tau$ et $F(s_0)$ un
repère orthonormé direct, a une unique solution globale, qui reste
un repère orthonormé direct~; (b) construire une courbe
birégulière paramétrée par \omterm{def:b2:curves:length}{longueur d'arc} de \omterm{thm:b2:curves:frenet2d}{courbure} $\kappa$ et
de \omterm{thm:b2:curves:frenet3d}{torsion} $\tau$~;
(c) prouver l'unicité à une isométrie directe de $\R^3$ près.
\end{exercise}

\section{Problème~: enveloppes --- l'astroïde, deux développées et
une caustique}

\begin{omfigure}
\begin{tikzpicture}[scale=2.2]
  \draw[->] (-1.15,0) -- (1.25,0) node[below] {$x$};
  \draw[->] (0,-1.15) -- (0,1.25) node[left] {$y$};
  \draw[omDef] (1,0) -- (0,0.0) (0.966,0) -- (0,0.259)
    (0.866,0) -- (0,0.5) (0.707,0) -- (0,0.707)
    (0.5,0) -- (0,0.866) (0.259,0) -- (0,0.966);
  \draw[omThm, very thick]
    plot[domain=0:360, samples=120]
    ({cos(\x)*cos(\x)*cos(\x)}, {sin(\x)*sin(\x)*sin(\x)});
\end{tikzpicture}

{\small Une échelle de \omterm{def:b2:curves:length}{longueur} $1$ glissant le long d'un mur
(positions en bleu) ne traverse jamais l'\omterm{pb:b2:curves:1}{astroïde} $x^{2/3} +
y^{2/3} = 1$
(en rouge)~: l'\omterm{pb:b2:curves:1}{astroïde} est l'\emph{\omterm{pb:b2:curves:1}{enveloppe}} de la famille de
segments, \omterm{def:b2:curves:arc}{tangente} à chacun d'eux.}
\end{omfigure}

\begin{problem}[{Problème de week-end --- la machine à enveloppes
et quatre courbes classiques}]\label{pb:b2:curves:1}
Une famille de droites à un paramètre échoue en général à
recouvrir le plan \omterm{def:b2:funcseq:def}{uniformément}~: les droites s'accumulent le long
d'une courbe \omterm{def:b2:curves:arc}{tangente} à toutes, leur \emph{enveloppe}\index{enveloppe}.
Les rayons lumineux
rendent les \omterm{pb:b2:curves:1}{enveloppes} visibles sous forme de \emph{caustiques}
--- la courbe brillante à rebroussement dans une tasse de café. Ce
problème construit la machine à \omterm{pb:b2:curves:1}{enveloppes} générale, puis la fait
tourner quatre fois~: l'échelle qui glisse (\omterm{pb:b2:curves:1}{astroïde}), les
normales de la parabole et de la cycloïde (développées, avec le
pendule de Huygens à la fin), et la caustique de la tasse de café
(\omterm{pb:b2:curves:1}{néphroïde}). Partout, $D_t$ désigne la droite d'équation
$a(t)\,x + b(t)\,y = c(t)$, où $a, b, c$
sont des fonctions $\mathcal C^2$ avec $(a(t), b(t)) \neq (0,0)$,
et $\Delta(t) = a(t)b'(t) - a'(t)b(t)$.

\medskip
\noindent\textbf{Partie I --- La machine à \omterm{pb:b2:curves:1}{enveloppes}.}
\begin{enumerate}
  \item Supposons $\Delta(t) \neq 0$. Montrer que le
        \emph{système caractéristique}
        \[
        \begin{cases} a(t)\,x + b(t)\,y = c(t)\\
        a'(t)\,x + b'(t)\,y = c'(t)\end{cases}
        \]
        a une unique solution $E(t) = (x(t), y(t))$, donnée par
        $x = \dfrac{cb' - c'b}{\Delta}$, $y = \dfrac{ac' -
        a'c}{\Delta}$.
  \item Supposons de plus que $E$ est $\mathcal C^1$ près de $t$
        avec $E'(t) \neq 0$. En dérivant la première
        équation du système, montrer $a(t)\,x'(t) +
        b(t)\,y'(t) = 0$, et en conclure que la courbe $E$
        passe par un point de $D_t$ \emph{avec la
        direction de} $D_t$~: la famille est \omterm{def:b2:curves:arc}{tangente} à $E$,
        qu'on appelle son \emph{\omterm{pb:b2:curves:1}{enveloppe}}.
  \item Vérification~: les \omterm{def:b2:curves:arc}{tangentes} de la parabole $y =
        x^2/2$ aux points $(t, t^2/2)$ sont $tx - y =
        t^2/2$. Vérifier que la machine à \omterm{pb:b2:curves:1}{enveloppes} redonne la
        parabole elle-même.
  \item (\omterm{pb:b2:curves:1}{Enveloppe} des normales) Soit $\gamma$ paramétré par
        \omterm{def:b2:curves:length}{longueur d'arc} avec $\kappa(s) \neq 0$. La normale en
        $\gamma(s)$ est $\{M : \langle M - \gamma(s), T(s)
        \rangle = 0\}$. Montrer que son système caractéristique
        force $\langle M - \gamma(s), N(s)\rangle =
        1/\kappa(s)$, d'où que le point caractéristique est
        le \omterm{def:b2:curves:curvature}{centre de courbure}~: \emph{l'\omterm{pb:b2:curves:1}{enveloppe} des
        normales est la développée}, retrouvant
        l'\cref{exo:b2:curves:6}. Vérifier $\Delta(s) =
        \kappa(s)$.
  \item Deux dégénérescences. Pour le faisceau $D_\theta : x\cos
        \theta + y\sin\theta = 0$, montrer que le
        point caractéristique est l'origine pour tout $\theta$
        (l'«~\omterm{pb:b2:curves:1}{enveloppe}~» se réduit à un point, et $E' = 0$~:
        la question 2 ne s'applique pas). Pour une famille de
        droites parallèles ($a, b$ constants), montrer $\Delta
        \equiv 0$ et
        que le système caractéristique est en général
        incompatible~: pas d'\omterm{pb:b2:curves:1}{enveloppe}.
\end{enumerate}

\medskip
\noindent\textbf{Partie II --- L'échelle qui glisse et
l'\omterm{pb:b2:curves:1}{astroïde}.} Un segment de \omterm{def:b2:curves:length}{longueur} $1$ glisse avec une extrémité
$P_t =
(\cos t, 0)$ sur le sol et l'autre $Q_t = (0, \sin t)$ sur
le mur, $t \in \intoo0{\pi/2}$.
\begin{enumerate}[resume]
  \item Montrer que la droite $(P_tQ_t)$ a pour équation $x\sin t +
        y\cos t = \sin t\cos t$, et que la machine à \omterm{pb:b2:curves:1}{enveloppes}
        donne le point caractéristique
        \[
        E(t) = (\cos^3 t,\ \sin^3 t) :
        \]
        l'\emph{astroïde}\index{astroïde}, d'équation implicite
        $x^{2/3} + y^{2/3} = 1$ (étendue aux
        autres quadrants par symétrie).
  \item Montrer que $E'(0) = 0$ et, en utilisant la classification
        locale (\cref{prop:b2:curves:local}), que
        l'\omterm{pb:b2:curves:1}{astroïde} a un point de rebroussement de première espèce
        en $(1, 0)$
        --- et de même en ses quatre points sur les axes.
  \item Calculer $\norm{E'(t)} = \tfrac32\abs{\sin 2t}$ et
        en déduire que la \omterm{def:b2:curves:length}{longueur} totale de l'\omterm{pb:b2:curves:1}{astroïde} est $6$.
  \item Où l'échelle touche-t-elle l'\omterm{pb:b2:curves:1}{astroïde}~? Montrer $E(t) =
        P_t + \sin^2 t\,(Q_t - P_t)$~: le point de contact
        divise l'échelle dans le rapport $\sin^2 t : \cos^2
        t$, la balayant d'une extrémité à l'autre à mesure que
        l'échelle glisse.
  \item Calculer l'aire délimitée par l'\omterm{pb:b2:curves:1}{astroïde}~: montrer que
        l'aire du premier quadrant est $3\int_0^{\pi/2}\sin^4
        t\cos^2 t\,\dd t$, évaluer l'intégrale par
        linéarisation ($\sin^2 2t = \tfrac{1 - \cos 4t}2$),
        et conclure que l'aire totale est $3\pi/8$.
\end{enumerate}

\medskip
\noindent\textbf{Partie III --- La développée de la parabole.}
Soit $\gamma(t) = (t, t^2/2)$.
\begin{enumerate}[resume]
  \item Montrer que la normale en $\gamma(t)$ a pour équation
        $x + t\,y = t + t^3/2$.
  \item Faire tourner la machine à \omterm{pb:b2:curves:1}{enveloppes}~: montrer que
        l'\omterm{pb:b2:curves:1}{enveloppe} des normales est
        \[
        E(t) = \Bigl(-t^3,\ 1 + \tfrac32 t^2\Bigr),
        \]
        d'équation implicite $x^2 = \tfrac8{27}(y - 1)^3$~:
        une parabole semi-cubique.
  \item Recouper avec la question 4~: calculer le \omterm{def:b2:curves:curvature}{centre de
        courbure} $\gamma(t) + \frac1{\kappa(t)}N(t)$ à partir de
        $\kappa(t) = (1 + t^2)^{-3/2}$
        (\cref{ex:b2:curves:kappaexamples}) et retrouver le
        même point.
  \item Montrer que la développée a un point de rebroussement de
        première espèce en
        $(0, 1)$, le \omterm{def:b2:curves:curvature}{centre de courbure} au sommet ---
        le point où $\kappa$ est extrémale, comme le prédit
        la formule $c' = -\frac{\kappa'}{\kappa^2}N$ de
        l'\cref{exo:b2:curves:6}.
  \item Combien de normales de la parabole passent par un
        point donné $(x_0, y_0)$~? Montrer que la réponse est
        gouvernée par la cubique $\tfrac{t^3}2 + (1 - y_0)\,t - x_0 = 0$~;
        traiter complètement le cas de l'axe $x_0 = 0$ (une
        normale pour $y_0 < 1$, trois pour $y_0 > 1$), et
        interpréter la développée comme
        la courbe de transition.
\end{enumerate}

\medskip
\noindent\textbf{Partie IV --- La caustique de la tasse de café.}
Des rayons parallèles de direction $(1, 0)$ frappent l'intérieur
du cercle miroir $x^2 + y^2 = 1$~; le rayon atteignant $P_\theta =
(\cos\theta, \sin\theta)$ se réfléchit selon la loi de la
réflexion.
\begin{enumerate}[resume]
  \item À partir de la symétrie miroir par rapport à la normale
        (le rayon),
        justifier que la direction réfléchie est $v = u -
        2\langle u, n\rangle n$ avec $u = (1,0)$, $n =
        (\cos\theta, \sin\theta)$, et calculer $v =
        -(\cos2\theta, \sin2\theta)$.
  \item Montrer que le rayon réfléchi est porté par la droite
        \[
        x\sin 2\theta - y\cos 2\theta = \sin\theta .
        \]
  \item Faire tourner la machine à \omterm{pb:b2:curves:1}{enveloppes} ($\Delta = 2$)~:
        montrer que la caustique est
        \[
        E(\theta) = \Bigl(\tfrac{3\cos\theta -
        \cos3\theta}4,\ \tfrac{3\sin\theta -
        \sin3\theta}4\Bigr),
        \]
        la \emph{néphroïde}\index{néphroïde}.
  \item Calculer $E'(\theta) = \tfrac32\sin\theta\,
        (\cos2\theta, \sin2\theta)$~; vérifier que la direction
        \omterm{def:b2:curves:arc}{tangente} est la direction du rayon réfléchi (question
        16), localiser les deux points de rebroussement
        $(\pm\tfrac12, 0)$, et
        montrer que le rayon réfléchi coupe l'axe $y = 0$
        en $x = \frac1{2\cos\theta}$ --- de sorte que les rayons
        presque axiaux se focalisent en $x = \tfrac12$~: la
        distance focale $R/2$ d'un miroir de rayon $R$.
  \item Montrer que la \omterm{pb:b2:curves:1}{néphroïde} a pour \omterm{def:b2:curves:length}{longueur} totale $6$ et que
        près de $\theta = 0$,
        \[
        E(\theta) - \bigl(\tfrac12, 0\bigr) =
        \bigl(\tfrac34\theta^2 + o(\theta^2),\ \theta^3 +
        o(\theta^3)\bigr) :
        \]
        un point de rebroussement de première espèce, pointant le
        long de l'axe.
  \item Expliquer en un paragraphe pourquoi la caustique est
        brillante~: par tout point juste à l'extérieur de la
        caustique passent deux
        rayons réfléchis, par tout point situé dessus les rayons
        sont «~infiniment concentrés~» (l'application $(\theta,
        \text{distance le long du rayon}) \mapsto \R^2$ a un
        point critique exactement sur l'\omterm{pb:b2:curves:1}{enveloppe}).
\end{enumerate}

\medskip
\noindent\textbf{Partie V --- Huygens~: la cycloïde est sa propre
développée.} Soit $\gamma(t) = (t - \sin t,\ 1 - \cos t)$, $t \in
\intoo0{2\pi}$, une arche de la cycloïde.
\begin{enumerate}[resume]
  \item Calculer $\kappa(t) = -\dfrac1{4\sin(t/2)}$ et le
        \omterm{def:b2:curves:curvature}{centre de courbure}~; montrer que la développée est
        \[
        c(t) = (t + \sin t,\ \cos t - 1),
        \]
        et que la substitution $t = u + \pi$ la présente
        comme la cycloïde originale translatée de $(\pi, -2)$~:
        \emph{la développée d'une cycloïde est une cycloïde
        congruente} (Huygens).
  \item Vérifier que le \omterm{def:b2:curves:curvature}{rayon de courbure} au sommet $t =
        \pi$ vaut $4$, soit la moitié de la \omterm{def:b2:curves:length}{longueur} $8$ d'une
        arche (\cref{exo:b2:curves:1})~; localiser le point de
        rebroussement de la développée directement sous le sommet,
        à distance $4$.
  \item (La propriété du fil) Soit $\gamma$ paramétré par \omterm{def:b2:curves:length}{longueur
        d'arc} avec
        $\kappa > 0$, $\kappa$ de classe $\mathcal C^1$ et
        $R = 1/\kappa$ strictement monotone. En utilisant $c' = R'N$,
        montrer que la \omterm{def:b2:curves:length}{longueur d'arc} de la développée entre
        $c(s_0)$ et $c(s_1)$ est $\abs{R(s_1) - R(s_0)}$.
        Interpréter~: un fil tendu déroulé de la développée,
        de \omterm{def:b2:curves:length}{longueur} $R(s_0)$ au départ, voit son extrémité libre
        tracer la courbe originale --- de sorte qu'un pendule
        oscillant entre deux joues cycloïdales de l'horloge de
        Huygens décrit une cycloïde.
  \item Synthèse. La machine de la partie I a produit
        l'\omterm{pb:b2:curves:1}{astroïde}, une parabole semi-cubique, une \omterm{pb:b2:curves:1}{néphroïde} et
        une cycloïde. Pour chacune des quatre familles, dire en une
        phrase où les hypothèses $\Delta \neq 0$ et
        $E' \neq 0$ étaient vérifiées ou en défaut, et quel
        événement géométrique (rebroussement, foyer,
        dégénérescence) chaque défaut de $E' \neq 0$
        a signalé. Où les extrema de \omterm{thm:b2:curves:frenet2d}{courbure} doivent-ils
        apparaître sur l'\omterm{pb:b2:curves:1}{enveloppe} des normales, et pourquoi~?
\end{enumerate}
\end{problem}
